\chapter{Subvariedades de \texorpdfstring{$\R^n$}{Rn}}
\label{ch:b3:submanifolds}

Esferas, toros, grupos de rotaciones: los hábitats naturales de la
geometría y la mecánica no son espacios vectoriales, sino
\emph{conjuntos curvos que de cerca parecen planos}. Este capítulo da a
esa frase un sentido preciso --- las subvariedades de $\R^n$ --- y el
cálculo para trabajar sobre ellas. El cimiento es el teorema de la
función inversa, demostrado aquí con el punto fijo de Banach; todo lo
demás es cambio de coordenadas: las cuatro descripciones equivalentes de
una subvariedad (enderezamiento local, conjuntos de nivel, grafos,
parametrizaciones), los \omterm{def:b3:submanifolds:tangent}{espacios
tangentes} y la optimización con restricciones por
\omterm{thm:b3:submanifolds:lagrange}{multiplicadores de Lagrange} ---
que, como demostración de despedida, vuelve a demostrar el teorema
espectral para matrices simétricas en tres líneas de geometría. El
problema de fin de semana construye el grupo de rotaciones $SO(3)$ y su
recubrimiento doble cuaterniónico: el álgebra (el $Q_8$
del~\cref{ch:b3:groups}, ya crecido) encontrándose con la geometría.

\section{El teorema de la función inversa}

\begin{theorem}[Teorema de la función inversa]
\label{thm:b3:submanifolds:ift}
Sean $U \subseteq \R^n$ abierto, $f \colon U \to \R^n$ de clase
$\mathcal C^1$ y $a \in U$ con $Df(a)$ invertible. Entonces existen
abiertos $V \ni a$, $W \ni f(a)$
tales que $f \colon V \to W$ es una biyección de inversa $\mathcal C^1$,
y\index{teorema de la función inversa}
\[
D(f^{-1})(y) = \bigl(Df(f^{-1}(y))\bigr)^{-1}
\qquad (y \in W).
\]
Si $f$ es $\mathcal C^k$, también lo es $f^{-1}$.
\end{theorem}

\begin{proof}
Normalicemos: sustituyendo $f$ por $x \mapsto Df(a)^{-1}\bigl(f(a +
x) - f(a)\bigr)$, podemos suponer $a = 0$, $f(0) = 0$, $Df(0) =
I$ (el enunciado general se sigue componiendo con las biyecciones
afines). Escríbase $f(x) = x + g(x)$: $Dg(0) = 0$ y, por
\omterm{def:b3:topology:continuity}{continuidad} de $Dg$, elíjase
$r > 0$ con $\vertiii{Dg(x)} \leq \frac12$ en $\bar B(0, r)$; la
desigualdad del valor medio da allí $\norm{g(x) - g(x')} \leq \frac12
\norm{x - x'}$.

\emph{Biyectividad sobre un \omterm{def:b3:topology:topology}{entorno}.}
Para $y \in B(0,
\frac r2)$, resolver $f(x) = y$ significa hallar un punto fijo de
$\Phi_y(x) = y - g(x)$; $\Phi_y$ lleva $\bar B(0, r)$ en sí mismo
($\norm{\Phi_y(x)} \leq \norm y + \frac12\norm x \leq
r$) y es $\frac12$-lipschitziana: Banach
(el~\cref{thm:b3:complete:banach}) da una única solución
$x = \varphi(y) \in \bar B(0,r)$. Además, $f$ es inyectiva en
$\bar B(0,r)$:
\[
\norm{f(x) - f(x')} \geq \norm{x - x'} - \norm{g(x) - g(x')}
\geq \tfrac12\norm{x - x'} .
\tag{$*$}
\]
Póngase $W = B(0, \frac r2)$ y $V = f^{-1}(W)\cap B(0, r)$: abierto
(\omterm{def:b3:topology:continuity}{continuidad}), con
$f \colon V \to W$ biyectiva.

\emph{\omterm{def:b3:topology:continuity}{Continuidad} y
diferenciabilidad de la inversa.} $(*)$ dice que $\varphi = f^{-1}$ es
$2$-lipschitziana. Fíjese $y_0 = f(x_0)
\in W$; la invertibilidad de $A = Df(x_0)$ (su distancia a $I$ es
$\leq \frac12$: Neumann, la~\cref{prop:b3:banach:neumann}) y la
diferenciabilidad de $f$ dan, para $y = f(x)$ próximo a $y_0$:
\[
\varphi(y) - \varphi(y_0) - A^{-1}(y - y_0)
= -A^{-1}\bigl(f(x) - f(x_0) - A(x - x_0)\bigr)
= -A^{-1}\,o(\norm{x - x_0}) = o(\norm{y - y_0}),
\]
usando $(*)$ para convertir $\norm{x - x_0} \leq 2\norm{y - y_0}$:
$\varphi$ es diferenciable en $y_0$ con la diferencial inversa.
\omterm{def:b3:topology:continuity}{Continuidad} de
$y \mapsto D\varphi(y) =
Df(\varphi(y))^{-1}$: composición de
aplicaciones \omterm{def:b3:topology:continuity}{continuas} (la
inversión es \omterm{def:b3:topology:continuity}{continua},
\cref{prop:b3:banach:neumann}): $\varphi \in \mathcal C^1$;
y reiterando la misma fórmula se obtiene $\mathcal C^k$.
\end{proof}

\begin{theorem}[Teorema de la función implícita]
\label{thm:b3:submanifolds:implicit}
Sea $F \colon U \subseteq \R^p\times\R^q \to \R^q$ de clase
$\mathcal C^1$ cerca de $(a, b)$, con $F(a,b) = 0$, y supóngase que la
diferencial parcial $D_yF(a,b) \in \mathcal L(\R^q)$ es invertible.
Entonces existen \omterm{def:b3:topology:topology}{entornos} $A \ni a$,
$B \ni b$ y una aplicación $\mathcal C^1$ $\psi \colon A \to B$ con
\[
\bigl\{(x, y)\in A\times B : F(x,y) = 0\bigr\}
= \{(x, \psi(x)) : x \in A\},
\]
y $D\psi(x) = -D_yF(x, \psi(x))^{-1}\,D_xF(x,
\psi(x))$.\index{teorema de la función implícita}
\end{theorem}

\begin{proof}
Aplíquese el~\cref{thm:b3:submanifolds:ift} a $\Theta(x, y) = (x,
F(x,y))$: su diferencial en $(a,b)$, triangular por bloques con bloques
diagonales invertibles $I$ y $D_yF$, es invertible. La inversa local
tiene la forma $\Theta^{-1}(x, z) = (x, h(x,
z))$; póngase $\psi(x) = h(x, 0)$: entonces, localmente, $F(x, y) = 0$
si y solo si $\Theta(x,y) = (x, 0)$ si y solo si $y = \psi(x)$. La
fórmula: derívese $F(x, \psi(x)) = 0$ por la regla de la cadena.
\end{proof}

\section{Subvariedades: cuatro definiciones}

\begin{theorem}[Caracterizaciones equivalentes]
\label{thm:b3:submanifolds:characterizations}
Sean $M \subseteq \R^n$, $d \in \{0, \dots, n\}$ y $k \geq
1$. Son equivalentes, para cada punto $a \in M$ (y $M$ es una
\emph{subvariedad de dimensión $d$ y clase
$\mathcal C^k$}\index{subvariedad} si se cumplen en todo $a
\in M$):
\begin{enumerate}
  \item (\emph{Enderezamiento}) Existe un difeomorfismo $\mathcal C^k$
  $\Phi$ de un abierto $\Omega \ni a$ sobre un abierto
  $\Omega' \subseteq \R^n$ con
        \[
        \Phi(M\cap\Omega) = \Omega' \cap
        \bigl(\R^d\times\{0\}\bigr).
        \]
  \item (\emph{Conjunto de nivel}) Existe una \emph{sumersión}
  $\mathcal C^k$ $F \colon \Omega \to \R^{n-d}$ (es decir, con $DF(x)$
  sobreyectiva) en un abierto $\Omega \ni a$ con
  $M\cap\Omega = F^{-1}(0)$. \item (\emph{Grafo}) Salvo permutación de
  coordenadas, $M$ es localmente el grafo de una aplicación
  $\mathcal C^k$ $\psi
        \colon A \subseteq \R^d \to \R^{n-d}$. \item
  (\emph{Parametrización}) Existe una \emph{inmersión} $\mathcal C^k$
  $\varphi \colon A \subseteq \R^d \to
        \R^n$ ($D\varphi(u)$ inyectiva) con $A$ abierto, siendo
  $\varphi$ un \omterm{def:b3:topology:continuity}{homeomorfismo} de $A$
  sobre $M \cap
        \Omega$ para cierto abierto $\Omega \ni a$.
\end{enumerate}
\end{theorem}

\begin{proof}
(1)$\Rightarrow$(2): $F = (\Phi_{d+1}, \dots, \Phi_n)$ (últimas
coordenadas de $\Phi$): una sumersión ($D\Phi$ invertible).
(2)$\Rightarrow$(3): $DF(a)$ sobreyectiva: algún menor $q \times q$ del
jacobiano es invertible ($q = n - d$); tras permutar coordenadas,
$D_yF(a)$ es invertible, y el teorema de la función implícita
(el~\cref{thm:b3:submanifolds:implicit}) expresa $M$ localmente como un
grafo $y = \psi(x)$. (3)$\Rightarrow$(4): $\varphi(x) = (x, \psi(x))$:
una inmersión (diferencial $\bigl(\begin{smallmatrix}I\\
D\psi\end{smallmatrix}\bigr)$ inyectiva) y un
\omterm{def:b3:topology:continuity}{homeomorfismo} sobre el grafo
(inversa: la proyección, \omterm{def:b3:topology:continuity}{continua}).
(4)$\Rightarrow$(1): sea $\varphi(u_0) = a$;
complétese
$\operatorname{im}D\varphi(u_0)$ con un suplementario $E$
($\dim E = n - d$) y defínase $\Theta(u, v) = \varphi(u) + v$ en
$A\times E$: $D\Theta(u_0, 0)$ es biyectiva (su imagen contiene
$\operatorname{im}D\varphi(u_0)$ y $E$), de modo que $\Theta$ es un
difeomorfismo local (el~\cref{thm:b3:submanifolds:ift}); su inversa
$\Phi$ endereza: cerca de $a$, los puntos de $M$ son exactamente los
$\varphi(u) = \Theta(u, 0)$ --- para ello, la hipótesis de
\omterm{def:b3:topology:continuity}{homeomorfismo} en (4) garantiza que
$M\cap\Omega$, para $\Omega$ pequeño, no contiene otras hojas (hay que
excluir $\varphi(u') + v = m \in
M$ próximo a $a$ con $v \neq 0$ pequeño: $m = \varphi(u'')$ para cierto
$u''$ próximo a $u_0$ por la propiedad de
\omterm{def:b3:topology:continuity}{homeomorfismo}, y la inyectividad
local de $\Theta$ fuerza $v = 0$). Entonces
$\Phi(M\cap\Omega) = (A\times\{0\})
\cap\Phi(\Omega)$, restringiendo si hace falta.
\end{proof}

\begin{example}\label{ex:b3:submanifolds:examples}
La esfera $S^{n-1} = \{\norm x^2 = 1\}$: conjunto de nivel de la
sumersión $F(x) = \norm x_2^2 - 1$ en
$\R^n\setminus\{0\}$ ($DF(x) = 2x^{\mathsf T} \ne 0$):
una subvariedad $\mathcal C^\infty$ de dimensión $n - 1$. El toro en
$\R^3$: conjunto de nivel de $\bigl(\sqrt{x^2 + y^2} -
R\bigr)^2 + z^2 - r^2$ ($0 < r < R$). El cono $\{x^2 + y^2 =
z^2\}$ \emph{no} es una subvariedad en $0$
(el~\cref{exo:b3:submanifolds:1}). Grupos de matrices: $SL_n$ y $O_n$
son subvariedades de $M_n(\R)$
(el~\cref{exo:b3:submanifolds:5,exo:b3:submanifolds:6}) --- el punto de
partida de la teoría de Lie.
\end{example}

\section{Espacios tangentes}

\begin{definition}\label{def:b3:submanifolds:tangent}
Sean $M$ una $d$-subvariedad y $a \in M$. El \emph{espacio
tangente}\index{espacio tangente} $T_aM$ es el conjunto de los vectores
velocidad $\gamma'(0)$ de las curvas $\mathcal C^1$ $\gamma \colon
\intoo{-\varepsilon}\varepsilon \to M$ con $\gamma(0) = a$.
\end{definition}

\begin{proposition}\label{prop:b3:submanifolds:tangent}
$T_aM$ es un subespacio vectorial de $\R^n$ de dimensión $d$, y:
\begin{enumerate}
  \item si $M = F^{-1}(0)$ localmente con $F$ una sumersión:
        $T_aM = \ker DF(a)$;
  \item si $M$ está parametrizada por la inmersión $\varphi$
        ($\varphi(u_0) = a$): $T_aM =
        \operatorname{im}D\varphi(u_0)$.
\end{enumerate}
\end{proposition}

\begin{proof}
Las curvas de $M$ cumplen $F(\gamma(t)) = 0$; la regla de la cadena en
$0$ da $DF(a)\gamma'(0) = 0$: $T_aM \subseteq \ker DF(a)$.
Recíprocamente, el enderezamiento
(el~\cref{thm:b3:submanifolds:characterizations}(1)) transporta rectas
de $\R^d\times\{0\}$ a curvas de $M$: todo vector de un subespacio de
dimensión $d$ se realiza; comparando dimensiones
($\dim\ker DF(a) = n - (n - d) = d$) se fuerza la igualdad en (1), y el
mismo argumento de transporte da (2) ($D\varphi(u_0)$ aplicada a rectas
de $A$; de nuevo las dimensiones).
\end{proof}

\begin{theorem}[Multiplicadores de Lagrange]
\label{thm:b3:submanifolds:lagrange}
Sea $M = F^{-1}(0)$ con $F = (F_1, \dots, F_q) \colon \Omega
\to \R^q$ una sumersión $\mathcal C^1$, y sea $f \colon
\Omega \to \R$ de clase $\mathcal C^1$. Si la restricción
$f\restriction_M$ tiene un extremo local en $a \in M$, entonces existen
números reales únicos $\lambda_1, \dots, \lambda_q$ (los
\emph{\omterm{thm:b3:submanifolds:lagrange}{multiplicadores de
Lagrange}}) con\index{multiplicadores de Lagrange}
\[
\nabla f(a) = \lambda_1\nabla F_1(a) + \dots +
\lambda_q\nabla F_q(a) .
\]
\end{theorem}

\begin{proof}
Para toda curva $\gamma$ de $M$ que pase por $a$: $t \mapsto
f(\gamma(t))$ tiene un extremo local en $0$, de modo que $0 =
\frac{\dd}{\dd t}f(\gamma(t))\big|_0 = \langle\nabla f(a),
\gamma'(0)\rangle$: $\nabla f(a) \perp T_aM = \ker DF(a)$
(la~\cref{prop:b3:submanifolds:tangent}). Ahora bien, $\ker DF(a)^\perp
= \operatorname{im}DF(a)^{\mathsf T}$ (la identidad de rango y
ortogonalidad de segundo año, o el~\cref{exo:b3:hilbert:8} en dimensión
finita: $(\ker T)^\perp = \operatorname{im}T^*$), que está generado por
los gradientes $\nabla F_i(a)$ --- independientes, al ser $DF(a)$
sobreyectiva: los multiplicadores existen y son únicos.
\end{proof}

\begin{example}[El teorema espectral, geométricamente]
\label{ex:b3:submanifolds:rayleigh}
Sea $A$ una matriz real simétrica $n\times n$ y maximícese
$f(x) = \langle Ax, x\rangle$ en la esfera $S^{n-1}$
(compacta: el máximo se alcanza, en
cierto $v_1$). Lagrange con $F(x) = \norm x^2 - 1$: $\nabla f = 2Ax$ y
$\nabla F =
2x$ dan $Av_1 = \lambda_1v_1$ --- un vector propio, con
$\lambda_1 = \max_{S^{n-1}}\langle Ax, x\rangle$. Restrínjase $A$ a
$v_1^\perp$ (invariante: $\langle Av, v_1\rangle =
\langle v, Av_1\rangle = \lambda_1\langle v, v_1\rangle = 0$) e itérese:
una base ortonormal de vectores propios. El teorema espectral de segundo
año, vuelto a demostrar por pura optimización --- y la sombra en
dimensión infinita del mismo argumento demostró
el~\cref{lem:b3:spectral:existence}.
\end{example}

\begin{method}\label{met:b3:submanifolds:method}
Para demostrar que un conjunto es una subvariedad: exhíbase localmente
como $F^{-1}(0)$ con $DF$ sobreyectiva \emph{sobre el conjunto} (la vía
más común), o como un grafo. Para calcular su dimensión y su
\omterm{def:b3:submanifolds:tangent}{espacio tangente}: $d = n - q$ y
$T_a = \ker DF(a)$. Para optimizar sobre él: Lagrange --- compruébese
siempre antes la \omterm{def:b3:topology:compact}{compacidad} (o la coercividad), de modo que exista un
extremo al que aplicar el teorema, y recuérdese que la ecuación de los
multiplicadores es solo \emph{necesaria}: recójanse todos los puntos
críticos y compárense después los valores. Para grupos de matrices,
derívense curvas en la identidad para identificar los
\omterm{def:b3:submanifolds:tangent}{espacios tangentes}.
\end{method}

\section{Ejercicios}

\begin{exercise}[$\star$]\label{exo:b3:submanifolds:1}
(a) Verifíquese que los siguientes son subvariedades $\mathcal C^\infty$
y dense sus dimensiones: $S^{n-1}$; el hiperboloide
$\{x^2 + y^2 - z^2 = 1\}$; el toro del
\cref{ex:b3:submanifolds:examples}.
(b) Demuéstrese que el cono $C = \{x^2 + y^2 = z^2\} \subseteq
\R^3$ no es una $2$-subvariedad en $0$: determínese el número de
\omterm{def:b3:topology:components}{componentes conexas} de
$\bigl(C\setminus\{0\}\bigr)\cap B(0,\varepsilon)$ y compárese con el
que forzaría un enderezamiento
(el~\cref{thm:b3:submanifolds:characterizations}(1)) para un plano menos
un punto.
\end{exercise}

\begin{exercise}[$\star$]\label{exo:b3:submanifolds:2}
Calcúlense los \omterm{def:b3:submanifolds:tangent}{espacios tangentes}:
(a) $T_aS^{n-1}$ para cualquier $a$ (respuesta: $a^\perp$); (b) el plano
tangente al toro del~\cref{ex:b3:submanifolds:examples} en un punto
arbitrario del ecuador exterior $\{z = 0,\ x^2 + y^2 = (R + r)^2\}$; (c)
la recta tangente a la hélice $\varphi(t) = (\cos t,
\sin t, t)$ en $\varphi(t_0)$, comprobando
\cref{prop:b3:submanifolds:tangent}(2).
\end{exercise}

\begin{exercise}[$\star\star$]\label{exo:b3:submanifolds:3}
Sea $f(x, y) = (x^2 - y^2,\ 2xy)$ (es decir, $z \mapsto z^2$). (a) ¿En
qué puntos se aplica el~\cref{thm:b3:submanifolds:ift}? (b) Demuéstrese
que $f$ es localmente pero no globalmente invertible en
$\R^2\setminus\{0\}$, y exhíbanse explícitamente las dos inversas
locales definidas en un \omterm{def:b3:topology:topology}{entorno} de
$(1, 0)$ (las dos ramas de la raíz cuadrada). (c) Misma discusión para
la aplicación de \omterm{ex:b3:product:polar}{coordenadas polares} $(r,
\theta) \mapsto (r\cos\theta, r\sin\theta)$.
\end{exercise}

\begin{exercise}[$\star\star$]\label{exo:b3:submanifolds:4}
(Folium) Sean $F(x,y) = x^3 + y^3 - 3xy$ y $\mathcal C =
F^{-1}(0)$. (a) Demuéstrese que, cerca de todo punto de $\mathcal C$
distinto del origen, $\mathcal C$ es una $1$-subvariedad, localmente un
grafo en $x$ o en $y$ (¿cuál, dónde?). (b) Calcúlese la recta tangente
en $(\frac32, \frac32)$. (c) ¿Qué ocurre en el origen? (Se cruzan dos
ramas: exhíbanse dos curvas $\mathcal C^1$ de $\mathcal C$ que pasen por
$0$ con velocidades independientes, y conclúyase que no existe
enderezamiento.)
\end{exercise}

\begin{exercise}[$\star\star$]\label{exo:b3:submanifolds:5}
Sea $F(M) = M^{\mathsf T}M$ de $M_n(\R)$ en el espacio $S_n$ de las
matrices simétricas. (a) Demuéstrese que
$DF(M)(H) = M^{\mathsf T}H + H^{\mathsf T}M$ y que $DF(M)$ es
sobreyectiva sobre $S_n$ en todo $M \in O_n$ \emph{(dada $S \in S_n$,
pruébese con $H = \frac12MS$)}. (b) Conclúyase que $O_n = F^{-1}(I)$ es
una subvariedad $\mathcal C^\infty$
compacta de dimensión
$\frac{n(n-1)}2$, con $T_IO_n = \{H : H^{\mathsf T} = -H\}$, las
matrices antisimétricas. (c) Demuéstrese que $\eu^{tH} \in O_n$ para
toda $H$ antisimétrica: las direcciones tangentes se integran en curvas
dentro del grupo.
\end{exercise}

\begin{exercise}[$\star\star$]\label{exo:b3:submanifolds:6}
(a) Demuéstrese que $\det \colon M_n(\R) \to \R$ tiene diferencial
$D\det(M)(H) = \operatorname{tr}\bigl(\operatorname{com}(M)
^{\mathsf T}H\bigr)$, no nula en todo $M \in SL_n$. (b) Conclúyase que
$SL_n(\R)$ es una subvariedad de dimensión $n^2 - 1$ con
$T_ISL_n = \{H : \operatorname{tr}H = 0\}$. (c) ¿Es $GL_n(\R)$ una
subvariedad? ¿De qué dimensión?
\end{exercise}

\begin{exercise}[$\star\star$]\label{exo:b3:submanifolds:7}
Por \omterm{thm:b3:submanifolds:lagrange}{multiplicadores de Lagrange}:
(a) hállense los extremos de $f(x,y) = xy$ en la circunferencia $x^2 +
y^2 = 1$; (b) demuéstrese que, entre todos los vectores de probabilidad
$(p_1, \dots,
p_n)$ (positivos y de suma $1$), la entropía $-\sum
p_i\ln p_i$ se maximiza exactamente en la distribución uniforme; (c)
hállese el punto de la elipse $\{x^2/4 + y^2 = 1\}$ más próximo a
$(1, 0)$, y compruébese geométricamente la ecuación de los
multiplicadores (alineamiento normal).
\end{exercise}

\begin{exercise}[$\star\star\star$]\label{exo:b3:submanifolds:8}
Desarróllese por completo el~\cref{ex:b3:submanifolds:rayleigh}:
demuéstrese por inducción que una matriz real simétrica admite una base
ortonormal de vectores propios, siendo
$\lambda_1 \geq \dots \geq \lambda_n$ los sucesivos máximos con
restricción del cociente de Rayleigh. Dedúzcanse después las fórmulas de
Courant--Fischer del~\cref{exo:b3:spectral:8} en dimensión finita
directamente de esta construcción.
\end{exercise}

\begin{exercise}[$\star\star\star$]\label{exo:b3:submanifolds:9}
(Desigualdad de Hadamard) Para $M \in GL_n(\R)$ de columnas
$c_1, \dots, c_n$:
\[
\abs{\det M} \;\leq\; \prod_{i=1}^n\norm{c_i}_2 ,
\]
con igualdad si y solo si las columnas son ortogonales. \emph{(Redúzcase
a columnas de norma $1$ reescalando; maximícese $\det$ en el producto
\omterm{def:b3:topology:compact}{compacto} de esferas $(S^{n-1})^n$; en
un maximizador, Lagrange en cada columna por separado da
$\nabla_{c_i}\det =
\lambda_ic_i$, y $\nabla_{c_i}\det$ es la $i$-ésima columna de
$\operatorname{com}(M)$: dedúzcase que $M^{\mathsf T}M$ es diagonal,
luego $= I$ y por tanto $\det = \pm1$.)} Lectura geométrica: el volumen
de un paralelepípedo es a lo sumo el producto de las longitudes de sus
aristas.
\end{exercise}

\begin{exercise}[$\star\star$]\label{exo:b3:submanifolds:10}
¿Cerca de cuáles de sus puntos es la circunferencia $S^1$ un grafo $y =
\psi(x)$? ¿Un grafo $x = \chi(y)$? Verifíquese explícitamente la
caracterización por grafos
(el~\cref{thm:b3:submanifolds:characterizations}(3)) en $(1, 0)$, y
explíquese en una frase por qué \emph{alguna} permutación de coordenadas
basta siempre pero ninguna fija sirve siempre.
\end{exercise}

\begin{exercise}[$\star\star$]\label{exo:b3:submanifolds:11}
(El grupo ortogonal como subvariedad, cuantitativamente) (a) Demuéstrese
que $O_n = \{M : M^{\mathsf T}M = I\}$ es
\omterm{def:b3:topology:compact}{compacto}: acotado (cada columna es un
vector unitario, de modo que $\norm M \leq
\sqrt n$ para la norma matricial euclídea) y cerrado. (b) Demuéstrese
que su \omterm{def:b3:submanifolds:tangent}{espacio tangente} en $I$ es
el espacio de las matrices antisimétricas, de dimensión
$\frac{n(n-1)}2$, y en un $A \in O_n$ general:
$T_AO_n = \{AK : K^{\mathsf T} =
-K\}$. (c) Dedúzcase que la aplicación $t \mapsto A\exp(tK)$ es, para
cada $K$ antisimétrica, una curva en $O_n$ que pasa por $A$ con
velocidad $AK$ \emph{(verifíquese $\exp(tK) \in O_n$ usando
$\exp(X)^{\mathsf T} = \exp(X^{\mathsf T})$ y $\exp(-X)\exp(X) = I$)}:
todo vector tangente se realiza mediante una curva explícita, sin
necesidad alguna del teorema de la función implícita.
\end{exercise}

\begin{exercise}[$\star\star$]\label{exo:b3:submanifolds:12}
(Puntos críticos de la distancia) Sean $M \subseteq \R^n$ una
subvariedad y $p \notin M$. Demuéstrese que si $x_0 \in M$ minimiza la
distancia a $p$ (tal punto existe cuando $M$ es cerrada y no vacía ---
¿por qué?), entonces
\[
p - x_0 \;\perp\; T_{x_0}M
\]
\emph{(derívese $t \mapsto \norm{\gamma(t) - p}^2$ a lo largo de curvas
de $M$)}. Dedúzcase: el punto más próximo de una esfera está en el radio
que pasa por el \omterm{ex:b3:groups:actions}{centro}; y úsese la
condición para calcular la distancia de $p = (2, 0)$ a la parábola
$y = x^2$ (redúzcase a una cúbica y resuélvase numéricamente con tres
cifras).
\end{exercise}

\section{Problema: \texorpdfstring{$SO(3)$}{SO(3)} y los
cuaterniones}

\begin{problem}[{Problema de fin de semana --- rotaciones, el grupo
$S^3$ y el recubrimiento doble}]\label{pb:b3:submanifolds:1} Los
cuaterniones $\mathbb H = \{t + x\mathrm i + y\mathrm j +
z\mathrm k\}$ --- el álgebra cuyo grupo de unidades contiene el $Q_8$
del~\cref{pb:b3:groups:1} --- parametrizan las rotaciones
tridimensionales por partida doble: la aplicación «conjugar por un
cuaternión unitario» es un morfismo sobreyectivo $S^3 \to
SO(3)$ de núcleo $\{\pm1\}$. Lo construiremos todo. Recuérdese o
defínase: la multiplicación es $\R$-bilineal con
$\mathrm i^2 = \mathrm j^2 = \mathrm k^2 =
\mathrm{ijk} = -1$; el conjugado de $q = t + x\mathrm i +
y\mathrm j + z\mathrm k$ es $\bar q = t - x\mathrm i -
y\mathrm j - z\mathrm k$; $N(q) = q\bar q = t^2 + x^2 + y^2
+ z^2$.

\medskip
\noindent\textbf{Parte I --- El álgebra $\mathbb H$ y el grupo $S^3$.}
\begin{enumerate}
  \item Verificar que $\mathbb H$ es un álgebra asociativa sobre $\R$ de
  \omterm{ex:b3:groups:actions}{centro} $\R$, que $\overline{pq} =
        \bar q\,\bar p$ y que $N(pq) = N(p)N(q)$ \emph{(una vía limpia:
  represéntese $q$ como la matriz compleja $2\times2$
  $\bigl(\begin{smallmatrix}\alpha & \beta\\
        -\bar\beta & \bar\alpha\end{smallmatrix}\bigr)$, $q =
        \alpha + \beta\mathrm j$, y úsese $\det$)}. \item Dedúzcase que
  todo $q \neq 0$ es invertible ($q^{-1} =
        \bar q/N(q)$): $\mathbb H$ es un cuerpo (no conmutativo) y
  $S^3 = \{N(q) = 1\}$ es un grupo --- y una $3$-subvariedad
  compacta de $\R^4$
        (\cref{ex:b3:submanifolds:examples}).
\end{enumerate}

\medskip
\noindent\textbf{Parte II --- El morfismo de rotación.} Identifíquese
$\R^3$ con los \emph{cuaterniones puros} $P = \{x\mathrm i +
y\mathrm j + z\mathrm k\}$, y para $q \in S^3$ defínase
$\rho_q(v) = q\,v\,\bar q$.
\begin{enumerate}[resume]
  \item Demostrar que $\rho_q$ lleva $P$ en $P$ \emph{(los cuaterniones
  puros son los de $\bar v = -v$)}, es $\R$-lineal, conserva la norma, y
  que $\rho
        \colon q \mapsto \rho_q$ es un morfismo de grupos $S^3
        \to O(3)$. \item Calcular el núcleo: $\rho_q = \mathrm{id}$ si y
  solo si $q$ conmuta con $\mathrm i, \mathrm j, \mathrm k$ si y solo si
        $q \in \R\cap S^3 = \{\pm1\}$.
  \item Escríbase $q = \cos\frac\theta2 + \sin\frac\theta2\,u$ con
  $u \in P$, $N(u) = 1$ (¿por qué es siempre posible para $q \in S^3$?).
  Demostrar que $\rho_q$ fija $u$ y, en el plano $u^\perp\cap P$, actúa
  como la rotación de ángulo $\theta$ \emph{(calcúlese $\rho_q(w)$ para
  $w \perp u$ usando $uw = -wu$ para unidades puras ortogonales ---
  demuéstrese esta identidad a partir de la tabla de multiplicar, o de
  $uw + wu =
        -2\langle u, w\rangle$)}. \item Conclúyase:
  $\operatorname{im}\rho \subseteq SO(3)$ (cada $\rho_q$ es una rotación
  de eje y ángulo los calculados --- determinante $+1$ por
  \omterm{def:b3:topology:continuity}{continuidad} de $q
        \mapsto \det\rho_q$ en el
  \omterm{def:b3:topology:connected}{conexo} $S^3$, o directamente), y
  $\rho$ es \emph{sobreyectivo} sobre $SO(3)$: toda rotación de $\R^3$
  tiene un eje (demuéstrese: una matriz ortogonal real $3\times3$ con
  $\det = 1$ tiene el valor propio $1$ --- considérese el polinomio
  característico) y es, por tanto, algún $\rho_q$. Resumen:
        \[
        SO(3) \;\cong\; S^3/\{\pm 1\} .
        \]
\end{enumerate}

\medskip
\noindent\textbf{Parte III --- $SO(3)$ como subvariedad; Rodrigues.}
\begin{enumerate}[resume]
  \item Demostrar que $SO(3)$ es una subvariedad
  compacta de dimensión $3$ de
  $M_3(\R)$ con $T_ISO(3) =$ las matrices antisimétricas
  (el~\cref{exo:b3:submanifolds:5}; la condición sobre el determinante
  selecciona una unión de componentes). \item Para la matriz
  antisimétrica $A_u$ asociada a $u \in \R^3$ ($A_uv = u\wedge v$, el
  producto vectorial), demuéstrese la \emph{fórmula de Rodrigues}:
        \[
        \eu^{\theta A_u} = I + \sin\theta\,A_u + (1 -
        \cos\theta)\,A_u^2 \qquad (\norm u = 1)
        \]
        \emph{(a partir de $A_u^3 = -A_u$: pártase la serie exponencial
        según las potencias de $A_u$)}, e identifíquese como la rotación
        de eje $u$ y ángulo $\theta$. Dedúzcase que $\exp$ lleva las
        matrices antisimétricas \emph{sobre}
        $SO(3)$.
  \item Relációnense las dos parametrizaciones: demuéstrese que $t
        \mapsto \rho_{q(t)}$ con $q(t) = \cos\frac t2 +
        \sin\frac t2\,u$ es un grupo uniparamétrico de rotaciones cuya
  derivada en $t = 0$ es $A_u$ --- la exponencial cuaterniónica y la
  \omterm{thm:b3:ode:matrixexp}{exponencial de matrices} cuentan la
  misma historia a media velocidad y a velocidad \omterm{def:b3:complete:complete}{completa},
  respectivamente.
\end{enumerate}

\medskip
\noindent\textbf{Parte IV --- El recubrimiento doble, palpado.}
\begin{enumerate}[resume]
  \item Demostrar que el \omterm{def:b3:topology:pathconnected}{camino}
  $q(t) = \cos\frac t2 + \sin\frac
        t2\,\mathrm k$, $t \in \intcc0{2\pi}$, es un lazo en $SO(3)$ (su
  imagen $\rho_{q(t)}$ vuelve a la identidad) cuyo levantamiento
  cuaterniónico \emph{no} es un lazo: $q(2\pi) = -q(0)$. Continuando
  hasta $t = 4\pi$ se cierra el levantamiento. Explíquese en un párrafo
  breve qué significa esto: una rotación de $2\pi$ no se puede deshacer
  \omterm{def:b3:topology:continuity}{continuamente} mientras que una de
  $4\pi$ sí (el truco del cinturón), porque los lazos de $SO(3)$ se
  detectan en su recubrimiento doble $S^3$. \item Dedúzcase también el
  dividendo práctico: componer rotaciones = multiplicar cuaterniones
  ($4$ multiplicaciones de datos en lugar de $9$, sin deriva de la
  ortogonalidad) --- verifíquese con la composición de dos cuartos de
  vuelta alrededor de $\mathrm i$ y $\mathrm j$: calcúlense el eje y el
  ángulo del producto.
\end{enumerate}

\medskip
\noindent\textbf{Parte V --- La matriz explícita: Euler--Rodrigues.}
Escríbase $q = a + b\mathrm i + c\mathrm j +
d\mathrm k \in S^3$, de modo que $a^2 + b^2 + c^2 + d^2 = 1$.
\begin{enumerate}[resume]
  \item Calcúlese $\rho_q(\mathrm i)$ íntegramente a partir de la tabla
  de multiplicar; obténganse después $\rho_q(\mathrm
        j)$ y $\rho_q(\mathrm k)$ mediante la sustitución cíclica
  $\mathrm i \to \mathrm j \to \mathrm k
        \to \mathrm i$, $(b, c, d) \to (c, d, b)$ \emph{(justifíquese:
  la permutación cíclica $\mathrm i, \mathrm j,
        \mathrm k$ se extiende a un automorfismo de $\mathbb
        H$, porque las relaciones de definición son cíclicamente
  simétricas)}. Conclúyase que la matriz de $\rho_q$ en la base
  $(\mathrm i, \mathrm j, \mathrm k)$ es la \emph{matriz de
  Euler--Rodrigues}
        \[
        R_q = \begin{pmatrix}
        a^2 + b^2 - c^2 - d^2 & 2(bc - ad) & 2(bd + ac)\\
        2(bc + ad) & a^2 - b^2 + c^2 - d^2 & 2(cd - ab)\\
        2(bd - ac) & 2(cd + ab) & a^2 - b^2 - c^2 + d^2
        \end{pmatrix}.
        \]
  \item (Leer una rotación al revés) Demuéstrese que
        \[
        \operatorname{tr}R_q = 4a^2 - 1 = 1 + 2\cos\theta,
        \qquad
        \tfrac12\bigl(R_q - R_q^{\mathsf T}\bigr) =
        \sin\theta\,A_u,
        \]
        con la notación de las preguntas 5 y 8. Dedúzcase un algoritmo
        que recupere $\pm q$ a partir de una matriz de rotación $R$: el
        ángulo, de la traza; el eje, de la parte antisimétrica cuando
        $0 < \theta < \pi$; y, cuando $\theta = \pi$, demuéstrese y
        úsese la identidad $R
        + I = 2\,uu^{\mathsf T}$. \item Evalúese $R_q$ para el producto
        $q =
        \frac12(1 + \mathrm i + \mathrm j + \mathrm k)$ de la pregunta
        11: aparece una matriz de permutación. Identifíquese la rotación
        y concíliese con el eje y el ángulo hallados en la pregunta 11.
\end{enumerate}

\medskip
\noindent\textbf{Parte VI --- Dentro de $S^3$: $SU(2)$,
\omterm{ex:b3:groups:actions}{clases de conjugación}, exponenciales.}
\begin{enumerate}[resume]
  \item Demostrar que la
  \omterm{def:b3:representations:rep}{representación} matricial de la
  pregunta 1 (llámese $\Phi$) se restringe a un isomorfismo de grupos de
  $S^3$ sobre el \emph{grupo unitario especial}
        \[
        SU(2) = \bigl\{U \in M_2(\C) : U^*U = I,\ \det U =
        1\bigr\}
        \]
        \emph{(para la sobreyectividad, escríbanse las ecuaciones
        $U^{-1} = U^*$ y $\det U = 1$ para una matriz compleja
        $2\times2$ general)}. \item Demostrar que la parte real es un
        invariante de conjugación en $S^3$ ---
        $\operatorname{Re}(pq\bar p) =
        \operatorname{Re}q$ para todos $p \in S^3$ --- y,
        recíprocamente, que dos cuaterniones unitarios con la misma
        parte real son conjugados en $S^3$ \emph{(redúzcase a llevar un
        eje puro unitario sobre otro, cosa que proporciona la Parte
        II)}. Descríbanse geométricamente las
        \omterm{ex:b3:groups:actions}{clases de conjugación} de $S^3$;
        tradúzcase a $SU(2)$ (conjuntos de nivel de la traza); y
        proyéctese por $\rho$: dos rotaciones son conjugadas en $SO(3)$
        si y solo si tienen el mismo ángulo $\theta \in \intcc0\pi$.
        \item Defínase $\exp$ en $\mathbb H$ por la serie exponencial;
        compruébese la convergencia absoluta usando
        $\abs{pq} = \abs p\,\abs q$ para $\abs q =
        \sqrt{N(q)}$. Demuéstrese, para $u$ puro unitario y $\theta
        \in \R$,
        \[
        \exp(\theta u) = \cos\theta + \sin\theta\,u ,
        \]
        dedúzcase que $\exp$ lleva el hiperplano $P$ sobre $S^3$, y
        compruébese que $\rho_{\exp(su)} =
        \eu^{2sA_u}$: de nuevo el fenómeno del semiángulo de la pregunta
        9. \item Para cuaterniones puros $v, w$, demuéstrese la regla
        del producto $vw = -\langle v, w\rangle + v\wedge w$ y, por
        tanto, la identidad del
        \omterm{def:b3:groups:derived}{conmutador}
        $vw - wv = 2\,v\wedge w$; demuéstrese también
        $[A_v, A_w] = A_{v\wedge w}$ para las matrices de la pregunta 8.
        Conclúyase que la derivada de $\rho$ en $1$ a lo largo de las
        curvas $t \mapsto \exp(tv)$ es el isomorfismo lineal
        $v \mapsto 2A_v$ de $P$ sobre las matrices antisimétricas, y que
        transporta el \omterm{def:b3:groups:derived}{conmutador} de
        cuaterniones al \omterm{def:b3:groups:derived}{conmutador} de
        matrices.
\end{enumerate}

\medskip
\noindent\textbf{Parte VII --- Estructura global.}
\begin{enumerate}[resume]
  \item (No hay sección \omterm{def:b3:topology:continuity}{continua})
  Supóngase que $s \colon SO(3) \to
        S^3$ es \omterm{def:b3:topology:continuity}{continua} con
  $\rho \circ s =
        \operatorname{id}$. Para el lazo $R(t) =
        \rho_{q(t)}$ de la pregunta 10, póngase $\varepsilon(t) =
        s(R(t))\,q(t)^{-1}$ para $t \in \intcc0{2\pi}$. Demuéstrese que
  $\varepsilon$ es \omterm{def:b3:topology:continuity}{continua} con
  valores en $\{\pm1\}$, y dedúzcase una contradicción: no existe
  ninguna elección global \omterm{def:b3:topology:continuity}{continua}
  de un cuaternión unitario que represente cada rotación. \item (El
  modelo de la bola) Sean $\bar B \subseteq \R^3$ la bola cerrada de
  radio $\pi$ y $E(v) = \eu^{A_v}$, con $E(0) = I$. Demuéstrese que $E$
  lleva $\bar B$ sobre $SO(3)$, es inyectiva en la bola abierta y, en la
  esfera frontera, identifica exactamente los antípodas:
  $E(\pi u) = E(-\pi u) = 2uu^{\mathsf T} - I$, sin ninguna otra
  coincidencia. Así, $SO(3)$ es la bola con los puntos frontera
  antipodales pegados --- el espacio proyectivo $\mathbb{RP}^3$ --- y un
  diámetro se convierte en el lazo no contráctil de la pregunta 10.
  \item Demostrar que $\rho_p\rho_q\rho_p^{-1} = \rho_{pq\bar
        p}$; que las involuciones de $SO(3)$ (los $R \neq
        I$ con $R^2 = I$) son exactamente las medias vueltas $\rho_w$
  con $w$ un cuaternión puro unitario; y que el
  \omterm{ex:b3:groups:actions}{centro} de $SO(3)$ es trivial. \item
  Demostrar que toda rotación es producto de dos medias vueltas: para
  $q = \cos\frac\theta2 +
        \sin\frac\theta2\,u$, elíjase un $w \perp
        u$ puro unitario, compruébese que $w' = qw$ es de nuevo un
  cuaternión puro unitario y verifíquese $\rho_q =
        \rho_{w'}\rho_w$. ¿Dónde están los dos ejes y qué ángulo forman?
  \item Conclúyase el resumen
  \omterm{def:b3:topology:topology}{topológico}: $SO(3)$ es
  \omterm{def:b3:topology:compact}{compacto} y
  \omterm{def:b3:topology:pathconnected}{conexo por caminos} (dense dos
  demostraciones: imagen \omterm{def:b3:topology:continuity}{continua}
  de $S^3$ por $\rho$; imagen de $\exp$), mientras que $O(3)$ tiene
  exactamente dos \omterm{def:b3:topology:components}{componentes
  conexas}, cada una \omterm{def:b3:topology:continuity}{homeomorfa} a
  $SO(3)$.

  \item (Una composición, de tres maneras) Sean $R_1$ la rotación de
  $\frac\pi2$ alrededor del eje $z$ y $R_2$ la rotación de $\frac\pi2$
  alrededor del eje $x$. Calcúlense el eje y el ángulo de $R_2R_1$: (i)
  multiplicando las dos matrices $3\times3$ y usando la traza y la parte
  antisimétrica (Parte V); (ii) multiplicando los cuaterniones unitarios
  correspondientes $q_2q_1$. Compruébese que ambas respuestas coinciden:
  ángulo $\frac{2\pi}3$, eje $\frac1{\sqrt3}(1, -1, 1)$. \item (La
  \omterm{ex:b3:conformal:mobius}{transformación de Cayley}) Para $K$
  antisimétrica, demuéstrese que $I + K$ es invertible y
        \[
        C(K) = (I - K)(I + K)^{-1} \in SO(n),
        \]
        sin que $-1$ sea nunca valor propio de $C(K)$; demuéstrese que
        $K \mapsto C(K)$ es una biyección de las matrices antisimétricas
        sobre $\{R \in SO(n) : -1 \notin
        \operatorname{Sp}R\}$, con inversa $R \mapsto (I -
        R)(I + R)^{-1}$. (Una carta racional de $SO(n)$, compañera de la
        $\exp$ trascendente de
        \cref{exo:b3:submanifolds:11}.)
\end{enumerate}
\end{problem}
